% Options for packages loaded elsewhere
\PassOptionsToPackage{unicode}{hyperref}
\PassOptionsToPackage{hyphens}{url}
\PassOptionsToPackage{dvipsnames,svgnames,x11names}{xcolor}
%
\documentclass[
  letterpaper,
  DIV=11,
  numbers=noendperiod]{scrartcl}

\usepackage{amsmath,amssymb}
\usepackage{iftex}
\ifPDFTeX
  \usepackage[T1]{fontenc}
  \usepackage[utf8]{inputenc}
  \usepackage{textcomp} % provide euro and other symbols
\else % if luatex or xetex
  \usepackage{unicode-math}
  \defaultfontfeatures{Scale=MatchLowercase}
  \defaultfontfeatures[\rmfamily]{Ligatures=TeX,Scale=1}
\fi
\usepackage{lmodern}
\ifPDFTeX\else  
    % xetex/luatex font selection
\fi
% Use upquote if available, for straight quotes in verbatim environments
\IfFileExists{upquote.sty}{\usepackage{upquote}}{}
\IfFileExists{microtype.sty}{% use microtype if available
  \usepackage[]{microtype}
  \UseMicrotypeSet[protrusion]{basicmath} % disable protrusion for tt fonts
}{}
\makeatletter
\@ifundefined{KOMAClassName}{% if non-KOMA class
  \IfFileExists{parskip.sty}{%
    \usepackage{parskip}
  }{% else
    \setlength{\parindent}{0pt}
    \setlength{\parskip}{6pt plus 2pt minus 1pt}}
}{% if KOMA class
  \KOMAoptions{parskip=half}}
\makeatother
\usepackage{xcolor}
\setlength{\emergencystretch}{3em} % prevent overfull lines
\setcounter{secnumdepth}{-\maxdimen} % remove section numbering
% Make \paragraph and \subparagraph free-standing
\ifx\paragraph\undefined\else
  \let\oldparagraph\paragraph
  \renewcommand{\paragraph}[1]{\oldparagraph{#1}\mbox{}}
\fi
\ifx\subparagraph\undefined\else
  \let\oldsubparagraph\subparagraph
  \renewcommand{\subparagraph}[1]{\oldsubparagraph{#1}\mbox{}}
\fi

\usepackage{color}
\usepackage{fancyvrb}
\newcommand{\VerbBar}{|}
\newcommand{\VERB}{\Verb[commandchars=\\\{\}]}
\DefineVerbatimEnvironment{Highlighting}{Verbatim}{commandchars=\\\{\}}
% Add ',fontsize=\small' for more characters per line
\usepackage{framed}
\definecolor{shadecolor}{RGB}{241,243,245}
\newenvironment{Shaded}{\begin{snugshade}}{\end{snugshade}}
\newcommand{\AlertTok}[1]{\textcolor[rgb]{0.68,0.00,0.00}{#1}}
\newcommand{\AnnotationTok}[1]{\textcolor[rgb]{0.37,0.37,0.37}{#1}}
\newcommand{\AttributeTok}[1]{\textcolor[rgb]{0.40,0.45,0.13}{#1}}
\newcommand{\BaseNTok}[1]{\textcolor[rgb]{0.68,0.00,0.00}{#1}}
\newcommand{\BuiltInTok}[1]{\textcolor[rgb]{0.00,0.23,0.31}{#1}}
\newcommand{\CharTok}[1]{\textcolor[rgb]{0.13,0.47,0.30}{#1}}
\newcommand{\CommentTok}[1]{\textcolor[rgb]{0.37,0.37,0.37}{#1}}
\newcommand{\CommentVarTok}[1]{\textcolor[rgb]{0.37,0.37,0.37}{\textit{#1}}}
\newcommand{\ConstantTok}[1]{\textcolor[rgb]{0.56,0.35,0.01}{#1}}
\newcommand{\ControlFlowTok}[1]{\textcolor[rgb]{0.00,0.23,0.31}{#1}}
\newcommand{\DataTypeTok}[1]{\textcolor[rgb]{0.68,0.00,0.00}{#1}}
\newcommand{\DecValTok}[1]{\textcolor[rgb]{0.68,0.00,0.00}{#1}}
\newcommand{\DocumentationTok}[1]{\textcolor[rgb]{0.37,0.37,0.37}{\textit{#1}}}
\newcommand{\ErrorTok}[1]{\textcolor[rgb]{0.68,0.00,0.00}{#1}}
\newcommand{\ExtensionTok}[1]{\textcolor[rgb]{0.00,0.23,0.31}{#1}}
\newcommand{\FloatTok}[1]{\textcolor[rgb]{0.68,0.00,0.00}{#1}}
\newcommand{\FunctionTok}[1]{\textcolor[rgb]{0.28,0.35,0.67}{#1}}
\newcommand{\ImportTok}[1]{\textcolor[rgb]{0.00,0.46,0.62}{#1}}
\newcommand{\InformationTok}[1]{\textcolor[rgb]{0.37,0.37,0.37}{#1}}
\newcommand{\KeywordTok}[1]{\textcolor[rgb]{0.00,0.23,0.31}{#1}}
\newcommand{\NormalTok}[1]{\textcolor[rgb]{0.00,0.23,0.31}{#1}}
\newcommand{\OperatorTok}[1]{\textcolor[rgb]{0.37,0.37,0.37}{#1}}
\newcommand{\OtherTok}[1]{\textcolor[rgb]{0.00,0.23,0.31}{#1}}
\newcommand{\PreprocessorTok}[1]{\textcolor[rgb]{0.68,0.00,0.00}{#1}}
\newcommand{\RegionMarkerTok}[1]{\textcolor[rgb]{0.00,0.23,0.31}{#1}}
\newcommand{\SpecialCharTok}[1]{\textcolor[rgb]{0.37,0.37,0.37}{#1}}
\newcommand{\SpecialStringTok}[1]{\textcolor[rgb]{0.13,0.47,0.30}{#1}}
\newcommand{\StringTok}[1]{\textcolor[rgb]{0.13,0.47,0.30}{#1}}
\newcommand{\VariableTok}[1]{\textcolor[rgb]{0.07,0.07,0.07}{#1}}
\newcommand{\VerbatimStringTok}[1]{\textcolor[rgb]{0.13,0.47,0.30}{#1}}
\newcommand{\WarningTok}[1]{\textcolor[rgb]{0.37,0.37,0.37}{\textit{#1}}}

\providecommand{\tightlist}{%
  \setlength{\itemsep}{0pt}\setlength{\parskip}{0pt}}\usepackage{longtable,booktabs,array}
\usepackage{calc} % for calculating minipage widths
% Correct order of tables after \paragraph or \subparagraph
\usepackage{etoolbox}
\makeatletter
\patchcmd\longtable{\par}{\if@noskipsec\mbox{}\fi\par}{}{}
\makeatother
% Allow footnotes in longtable head/foot
\IfFileExists{footnotehyper.sty}{\usepackage{footnotehyper}}{\usepackage{footnote}}
\makesavenoteenv{longtable}
\usepackage{graphicx}
\makeatletter
\def\maxwidth{\ifdim\Gin@nat@width>\linewidth\linewidth\else\Gin@nat@width\fi}
\def\maxheight{\ifdim\Gin@nat@height>\textheight\textheight\else\Gin@nat@height\fi}
\makeatother
% Scale images if necessary, so that they will not overflow the page
% margins by default, and it is still possible to overwrite the defaults
% using explicit options in \includegraphics[width, height, ...]{}
\setkeys{Gin}{width=\maxwidth,height=\maxheight,keepaspectratio}
% Set default figure placement to htbp
\makeatletter
\def\fps@figure{htbp}
\makeatother

\usepackage{booktabs}
\usepackage{longtable}
\usepackage{array}
\usepackage{multirow}
\usepackage{wrapfig}
\usepackage{float}
\usepackage{colortbl}
\usepackage{pdflscape}
\usepackage{tabu}
\usepackage{threeparttable}
\usepackage{threeparttablex}
\usepackage[normalem]{ulem}
\usepackage{makecell}
\usepackage{xcolor}
\KOMAoption{captions}{tableheading}
\makeatletter
\@ifpackageloaded{tcolorbox}{}{\usepackage[skins,breakable]{tcolorbox}}
\@ifpackageloaded{fontawesome5}{}{\usepackage{fontawesome5}}
\definecolor{quarto-callout-color}{HTML}{909090}
\definecolor{quarto-callout-note-color}{HTML}{0758E5}
\definecolor{quarto-callout-important-color}{HTML}{CC1914}
\definecolor{quarto-callout-warning-color}{HTML}{EB9113}
\definecolor{quarto-callout-tip-color}{HTML}{00A047}
\definecolor{quarto-callout-caution-color}{HTML}{FC5300}
\definecolor{quarto-callout-color-frame}{HTML}{acacac}
\definecolor{quarto-callout-note-color-frame}{HTML}{4582ec}
\definecolor{quarto-callout-important-color-frame}{HTML}{d9534f}
\definecolor{quarto-callout-warning-color-frame}{HTML}{f0ad4e}
\definecolor{quarto-callout-tip-color-frame}{HTML}{02b875}
\definecolor{quarto-callout-caution-color-frame}{HTML}{fd7e14}
\makeatother
\makeatletter
\makeatother
\makeatletter
\makeatother
\makeatletter
\@ifpackageloaded{caption}{}{\usepackage{caption}}
\AtBeginDocument{%
\ifdefined\contentsname
  \renewcommand*\contentsname{Table of contents}
\else
  \newcommand\contentsname{Table of contents}
\fi
\ifdefined\listfigurename
  \renewcommand*\listfigurename{List of Figures}
\else
  \newcommand\listfigurename{List of Figures}
\fi
\ifdefined\listtablename
  \renewcommand*\listtablename{List of Tables}
\else
  \newcommand\listtablename{List of Tables}
\fi
\ifdefined\figurename
  \renewcommand*\figurename{Figure}
\else
  \newcommand\figurename{Figure}
\fi
\ifdefined\tablename
  \renewcommand*\tablename{Table}
\else
  \newcommand\tablename{Table}
\fi
}
\@ifpackageloaded{float}{}{\usepackage{float}}
\floatstyle{ruled}
\@ifundefined{c@chapter}{\newfloat{codelisting}{h}{lop}}{\newfloat{codelisting}{h}{lop}[chapter]}
\floatname{codelisting}{Listing}
\newcommand*\listoflistings{\listof{codelisting}{List of Listings}}
\makeatother
\makeatletter
\@ifpackageloaded{caption}{}{\usepackage{caption}}
\@ifpackageloaded{subcaption}{}{\usepackage{subcaption}}
\makeatother
\makeatletter
\@ifpackageloaded{tcolorbox}{}{\usepackage[skins,breakable]{tcolorbox}}
\makeatother
\makeatletter
\@ifundefined{shadecolor}{\definecolor{shadecolor}{rgb}{.97, .97, .97}}
\makeatother
\makeatletter
\makeatother
\makeatletter
\makeatother
\ifLuaTeX
  \usepackage{selnolig}  % disable illegal ligatures
\fi
\IfFileExists{bookmark.sty}{\usepackage{bookmark}}{\usepackage{hyperref}}
\IfFileExists{xurl.sty}{\usepackage{xurl}}{} % add URL line breaks if available
\urlstyle{same} % disable monospaced font for URLs
\hypersetup{
  pdftitle={Pretty tables with kableExtra},
  pdfauthor={Emily Malcolm-White},
  colorlinks=true,
  linkcolor={blue},
  filecolor={Maroon},
  citecolor={Blue},
  urlcolor={Blue},
  pdfcreator={LaTeX via pandoc}}

\title{Pretty tables with \texttt{kableExtra}}
\author{Emily Malcolm-White}
\date{}

\begin{document}
\maketitle
\ifdefined\Shaded\renewenvironment{Shaded}{\begin{tcolorbox}[frame hidden, interior hidden, borderline west={3pt}{0pt}{shadecolor}, sharp corners, boxrule=0pt, enhanced, breakable]}{\end{tcolorbox}}\fi

\includegraphics[width=0.3\textwidth,height=\textheight]{118_D_kable_files/mediabag/kableExtra_sm.png}

\begin{Shaded}
\begin{Highlighting}[]
\CommentTok{\#LOAD PACKAGES }
\FunctionTok{library}\NormalTok{(tidyverse)}
\end{Highlighting}
\end{Shaded}

\begin{Shaded}
\begin{Highlighting}[]
\CommentTok{\#LOAD DATA }
\FunctionTok{library}\NormalTok{(palmerpenguins)}
\FunctionTok{data}\NormalTok{(penguins)}

\CommentTok{\#CLEAN UP DATA}
\NormalTok{penguins }\OtherTok{\textless{}{-}}\NormalTok{ penguins }\SpecialCharTok{\%\textgreater{}\%}
  \FunctionTok{drop\_na}\NormalTok{()}
\CommentTok{\#sometimes this is appropriate. It\textquotesingle{}s questionable here... }
\end{Highlighting}
\end{Shaded}

\hypertarget{default-printing-style-of-a-table}{%
\section{default printing style of a
table}\label{default-printing-style-of-a-table}}

Let's consider our table from last class:

\begin{Shaded}
\begin{Highlighting}[]
\NormalTok{penguins }\SpecialCharTok{\%\textgreater{}\%}
  \FunctionTok{group\_by}\NormalTok{(species) }\SpecialCharTok{\%\textgreater{}\%}
  \FunctionTok{summarise}\NormalTok{(}\AttributeTok{average\_bill\_lenth =} \FunctionTok{mean}\NormalTok{(bill\_length\_mm), }
            \AttributeTok{average\_bill\_depth =} \FunctionTok{mean}\NormalTok{(bill\_depth\_mm))}
\end{Highlighting}
\end{Shaded}

\begin{verbatim}
# A tibble: 3 x 3
  species   average_bill_lenth average_bill_depth
  <fct>                  <dbl>              <dbl>
1 Adelie                  38.8               18.3
2 Chinstrap               48.8               18.4
3 Gentoo                  47.6               15.0
\end{verbatim}

When we knit this up it looks like of ugly\ldots{}

\hypertarget{using-kable-to-get-pretty-tables}{%
\section{\texorpdfstring{Using \texttt{kable} to get pretty
tables}{Using kable to get pretty tables}}\label{using-kable-to-get-pretty-tables}}

\begin{Shaded}
\begin{Highlighting}[]
\FunctionTok{library}\NormalTok{(kableExtra)}
\end{Highlighting}
\end{Shaded}

\begin{Shaded}
\begin{Highlighting}[]
\NormalTok{penguins }\SpecialCharTok{\%\textgreater{}\%}
  \FunctionTok{group\_by}\NormalTok{(species) }\SpecialCharTok{\%\textgreater{}\%}
  \FunctionTok{summarise}\NormalTok{(}\AttributeTok{average\_bill\_length =} \FunctionTok{mean}\NormalTok{(bill\_length\_mm), }
            \AttributeTok{average\_bill\_depth =} \FunctionTok{mean}\NormalTok{(bill\_depth\_mm)) }\SpecialCharTok{\%\textgreater{}\%} 
  \FunctionTok{kbl}\NormalTok{()}
\end{Highlighting}
\end{Shaded}

\begin{tabular}[t]{l|r|r}
\hline
species & average\_bill\_length & average\_bill\_depth\\
\hline
Adelie & 38.82397 & 18.34726\\
\hline
Chinstrap & 48.83382 & 18.42059\\
\hline
Gentoo & 47.56807 & 14.99664\\
\hline
\end{tabular}

\begin{Shaded}
\begin{Highlighting}[]
\CommentTok{\#OR }

\NormalTok{table1 }\OtherTok{\textless{}{-}}\NormalTok{ penguins }\SpecialCharTok{\%\textgreater{}\%}
  \FunctionTok{group\_by}\NormalTok{(species) }\SpecialCharTok{\%\textgreater{}\%}
  \FunctionTok{summarise}\NormalTok{(}\AttributeTok{average\_bill\_length =} \FunctionTok{mean}\NormalTok{(bill\_length\_mm), }
            \AttributeTok{average\_bill\_depth =} \FunctionTok{mean}\NormalTok{(bill\_depth\_mm))}

\FunctionTok{kbl}\NormalTok{(table1)}
\end{Highlighting}
\end{Shaded}

\begin{tabular}[t]{l|r|r}
\hline
species & average\_bill\_length & average\_bill\_depth\\
\hline
Adelie & 38.82397 & 18.34726\\
\hline
Chinstrap & 48.83382 & 18.42059\\
\hline
Gentoo & 47.56807 & 14.99664\\
\hline
\end{tabular}

\hypertarget{options-in-kable}{%
\section{Options in kable}\label{options-in-kable}}

We customize the content so it's displaying the information more
clearly.

\begin{Shaded}
\begin{Highlighting}[]
\NormalTok{table1 }\SpecialCharTok{\%\textgreater{}\%}
  \FunctionTok{kbl}\NormalTok{(}\AttributeTok{col.names =} \FunctionTok{c}\NormalTok{(}\StringTok{"Species"}\NormalTok{, }\StringTok{"Average Bill Length"}\NormalTok{, }\StringTok{"Average Bill Depth"}\NormalTok{), }
    \AttributeTok{caption =} \StringTok{"Average Bill Characteristics by Species"}\NormalTok{, }
    \AttributeTok{digits =} \DecValTok{2}\NormalTok{)}
\end{Highlighting}
\end{Shaded}

\begin{table}

\caption{Average Bill Characteristics by Species}
\centering
\begin{tabular}[t]{l|r|r}
\hline
Species & Average Bill Length & Average Bill Depth\\
\hline
Adelie & 38.82 & 18.35\\
\hline
Chinstrap & 48.83 & 18.42\\
\hline
Gentoo & 47.57 & 15.00\\
\hline
\end{tabular}
\end{table}

Better\ldots{}

\hypertarget{pretty-styling}{%
\section{pretty styling}\label{pretty-styling}}

\begin{Shaded}
\begin{Highlighting}[]
\NormalTok{table1 }\SpecialCharTok{\%\textgreater{}\%}
  \FunctionTok{kbl}\NormalTok{(}\AttributeTok{col.names =} \FunctionTok{c}\NormalTok{(}\StringTok{"Species"}\NormalTok{, }\StringTok{"Average Bill Length"}\NormalTok{, }\StringTok{"Average Bill Depth"}\NormalTok{), }
    \AttributeTok{caption =} \StringTok{"Average Bill Characteristics by Species"}\NormalTok{, }
    \AttributeTok{digits =} \DecValTok{2}\NormalTok{) }\SpecialCharTok{\%\textgreater{}\%}
  \FunctionTok{kable\_styling}\NormalTok{()}
\end{Highlighting}
\end{Shaded}

\begin{table}
\centering
\caption{Average Bill Characteristics by Species}
\centering
\begin{tabular}[t]{l|r|r}
\hline
Species & Average Bill Length & Average Bill Depth\\
\hline
Adelie & 38.82 & 18.35\\
\hline
Chinstrap & 48.83 & 18.42\\
\hline
Gentoo & 47.57 & 15.00\\
\hline
\end{tabular}
\end{table}

Many options for customizing the look of the tables -- more here:
\url{https://cran.r-project.org/web/packages/kableExtra/vignettes/awesome_table_in_html.html}

Let's try one\ldots{}

Let's make each row's color correspond to

\begin{Shaded}
\begin{Highlighting}[]
\NormalTok{table1 }\SpecialCharTok{\%\textgreater{}\%}
  \FunctionTok{kbl}\NormalTok{(}\AttributeTok{col.names =} \FunctionTok{c}\NormalTok{(}\StringTok{"Species"}\NormalTok{, }\StringTok{"Average Bill Length"}\NormalTok{, }\StringTok{"Average Bill Depth"}\NormalTok{), }
    \AttributeTok{caption =} \StringTok{"Average Bill Characteristics by Species"}\NormalTok{, }
    \AttributeTok{digits =} \DecValTok{2}\NormalTok{) }\SpecialCharTok{\%\textgreater{}\%}
  \FunctionTok{kable\_paper}\NormalTok{() }\SpecialCharTok{\%\textgreater{}\%}
  \FunctionTok{column\_spec}\NormalTok{(}\DecValTok{1}\NormalTok{, }\AttributeTok{bold=}\NormalTok{T) }\SpecialCharTok{\%\textgreater{}\%}
  \FunctionTok{row\_spec}\NormalTok{(}\DecValTok{2}\NormalTok{, }\AttributeTok{color =} \StringTok{"\#c85bcc"}\NormalTok{) }\SpecialCharTok{\%\textgreater{}\%}
  \FunctionTok{row\_spec}\NormalTok{(}\DecValTok{3}\NormalTok{, }\AttributeTok{color =} \StringTok{"\#067176"}\NormalTok{) }\SpecialCharTok{\%\textgreater{}\%}
  \FunctionTok{row\_spec}\NormalTok{(}\DecValTok{1}\NormalTok{, }\AttributeTok{color =} \StringTok{"\#ff7501"}\NormalTok{)}
\end{Highlighting}
\end{Shaded}

\begin{table}
\centering
\caption{Average Bill Characteristics by Species}
\centering
\begin{tabular}[t]{>{}l|r|r}
\hline
Species & Average Bill Length & Average Bill Depth\\
\hline
\textcolor[HTML]{ff7501}{\textbf{Adelie}} & \textcolor[HTML]{ff7501}{38.82} & \textcolor[HTML]{ff7501}{18.35}\\
\hline
\textcolor[HTML]{c85bcc}{\textbf{Chinstrap}} & \textcolor[HTML]{c85bcc}{48.83} & \textcolor[HTML]{c85bcc}{18.42}\\
\hline
\textcolor[HTML]{067176}{\textbf{Gentoo}} & \textcolor[HTML]{067176}{47.57} & \textcolor[HTML]{067176}{15.00}\\
\hline
\end{tabular}
\end{table}

RStudio hosts a table contest every year!

\begin{itemize}
\tightlist
\item
  Here is a link to this year's contest
  \url{https://www.rstudio.com/blog/rstudio-table-contest-2022/}
\item
  Here is a link to previous year's entries and winners to explore what
  is possible! \url{https://community.rstudio.com/c/table-gallery/64}
\end{itemize}

\begin{tcolorbox}[enhanced jigsaw, bottomrule=.15mm, rightrule=.15mm, colbacktitle=quarto-callout-tip-color!10!white, bottomtitle=1mm, left=2mm, toptitle=1mm, title=\textcolor{quarto-callout-tip-color}{\faLightbulb}\hspace{0.5em}{Tip}, toprule=.15mm, titlerule=0mm, breakable, arc=.35mm, colback=white, coltitle=black, opacitybacktitle=0.6, leftrule=.75mm, opacityback=0, colframe=quarto-callout-tip-color-frame]

As a general rule, you should have the content of the table as you'd
like it (the exact columns and rows you want) first and \emph{then} you
can make it pretty using the \texttt{kableExtra} package.

\end{tcolorbox}



\end{document}
